\documentclass[11pt]{article}

\usepackage[margin=1in]{geometry}
\usepackage{amsmath}
\usepackage{amssymb}
\usepackage{booktabs}
\usepackage{hyperref}
\usepackage{enumitem}

\hypersetup{
  colorlinks=true,
  linkcolor=blue,
  urlcolor=blue,
  citecolor=blue
}

\title{Starter Project Outline: Quantum Active-Space Astrochemistry for H + H$_2$CO on Ice}
\author{QDYNAMICS Project Notes}
\date{\today}

\begin{document}

\maketitle

\section{Core Idea}

The first chemically meaningful project should be simple enough to implement with the current
QDYNAMICS/Psi4 pipeline, but relevant enough to grow into the larger astrochemistry program.
The proposed starting reaction is hydrogen addition to formaldehyde,
\begin{align}
  \mathrm{H + H_2CO} &\rightarrow \mathrm{CH_3O}, \\
  \mathrm{H + H_2CO} &\rightarrow \mathrm{CH_2OH}.
\end{align}
This reaction is a natural first target because it is small, radical-like, relevant to grain-surface
methanol chemistry, and has competing product channels. The project can start in the gas phase,
then add water ice, then reduce the problem to an active-space Hamiltonian suitable for exact
diagonalization and ADAPT-VQE.

\section{Why This System}

\begin{itemize}[leftmargin=*]
  \item It is astrochemically motivated: hydrogenation of CO/H$_2$CO is central to methanol and complex organic molecule formation on dust-grain ice.
  \item It is small enough for a first active-space workflow.
  \item It has two channels, so the final observable can be a branching ratio rather than only a ground-state energy.
  \item It naturally introduces tunneling, because H motion controls the low-temperature rate.
  \item It can be embedded gradually in water: first H$_2$O, then (H$_2$O)$_4$, then (H$_2$O)$_8$.
\end{itemize}

\section{Computational Ladder}

The project should proceed through a ladder where each stage adds only one new complication.

\begin{center}
\begin{tabular}{lll}
\toprule
Stage & System & Purpose \\
\midrule
0 & H$_2$ & Verify Psi4, AO/MO integrals, spin-orbital indexing. \\
1 & H$_2$CO & Verify molecular integrals for the first real target molecule. \\
2 & H + H$_2$CO & Build reactant/product and stretched-bond geometries. \\
3 & H + H$_2$CO + H$_2$O & Add the first ice hydrogen-bond field. \\
4 & H + H$_2$CO + (H$_2$O)$_n$ & Move toward ice clusters, $n=4,8,16$. \\
5 & Active-space Hamiltonian & Downfold to CAS$(N_e,N_o)$ for quantum solvers. \\
\bottomrule
\end{tabular}
\end{center}

\section{Hamiltonian Target}

For a fixed nuclear geometry $R$, the electronic Hamiltonian in spin orbitals is
\begin{equation}
  \hat{H}_{\mathrm{el}}(R)
  =
  E_{\mathrm{nuc}}(R)
  +
  \sum_{pq} h_{pq}(R) a_p^\dagger a_q
  +
  \frac{1}{2}\sum_{pqrs} h_{pqrs}(R)
  a_p^\dagger a_q^\dagger a_s a_r .
\end{equation}
The current Psi4 pipeline is already moving in this direction by constructing
\begin{equation}
  h_{pq}
  \quad \mathrm{and} \quad
  h_{pqrs}
\end{equation}
from AO integrals and MO coefficients.

In chemist notation, the spatial two-electron integral is
\begin{equation}
  (ij|kl)
  =
  \int d\mathbf{r}_1 d\mathbf{r}_2\,
  \phi_i(\mathbf{r}_1)\phi_j(\mathbf{r}_1)
  \frac{1}{r_{12}}
  \phi_k(\mathbf{r}_2)\phi_l(\mathbf{r}_2).
\end{equation}
In spin-orbital form,
\begin{equation}
  \langle pq|rs\rangle
  =
  (ij|kl)\,
  \delta_{\sigma_p\sigma_q}\,
  \delta_{\sigma_r\sigma_s}.
\end{equation}
This is why the spin blocks in the code require the first pair to share spin and the second pair to share spin.

\section{First Reaction Geometry Set}

Before locating a full transition state, use a small controlled geometry set:
\begin{enumerate}[leftmargin=*]
  \item Reactant-like geometry: H far from the carbonyl carbon or oxygen.
  \item Midpoint geometry: H approaching either C or O along a chemically reasonable coordinate.
  \item Product-like geometry: CH$_3$O or CH$_2$OH radical geometry.
\end{enumerate}
This gives an initial reaction profile without yet needing a full nudged elastic band or transition-state search.

\section{Ice Extension}

The ice environment should be introduced gradually:
\begin{align}
  \mathrm{H + H_2CO}
  &\rightarrow
  \mathrm{products}, \\
  \mathrm{H + H_2CO + H_2O}
  &\rightarrow
  \mathrm{products + H_2O}, \\
  \mathrm{H + H_2CO + (H_2O)_n}
  &\rightarrow
  \mathrm{products + (H_2O)_n}.
\end{align}
The first target observable is the ice-induced barrier shift,
\begin{equation}
  \Delta\Delta E^\ddagger_{\mathrm{ice}}
  =
  \Delta E^\ddagger_{\mathrm{ice}}
  -
  \Delta E^\ddagger_{\mathrm{gas}}.
\end{equation}
If this quantity is negative, the ice lowers the barrier. If it is positive, the ice suppresses the channel.

\section{Active-Space Reduction}

The full molecular orbital space will become too large once water molecules are added. The project
therefore needs active-space reduction. The first chemically motivated active orbitals are:
\begin{itemize}[leftmargin=*]
  \item carbonyl $\pi$ orbital,
  \item carbonyl $\pi^\ast$ orbital,
  \item incoming H 1s orbital,
  \item radical/product SOMO,
  \item oxygen lone pair if hydrogen bonding or proton-coupled effects become important.
\end{itemize}
A plausible first active space is
\begin{equation}
  \mathrm{CAS}(5e,5o)
  \quad \mathrm{or} \quad
  \mathrm{CAS}(7e,7o),
\end{equation}
corresponding to 10 or 14 spin orbitals before qubit tapering.

\subsection{Active-Space Selection Strategy}

The first implementation can be manual and chemistry-driven. Later, it should become automatic:
\begin{enumerate}[leftmargin=*]
  \item Run a mean-field calculation for each geometry.
  \item Inspect orbital characters near the HOMO/LUMO window.
  \item Include orbitals involved in forming or breaking bonds.
  \item Add natural orbital occupation analysis when correlated references are available.
  \item Freeze inactive occupied orbitals and discard high-energy virtual orbitals.
\end{enumerate}

\section{Validation Targets}

For every geometry, validate the electronic data before running a quantum algorithm:
\begin{itemize}[leftmargin=*]
  \item Check electron count from the density matrix.
  \item Check one-electron energy.
  \item Check two-electron Hartree-Fock energy:
  \begin{equation}
    E_{2e}^{\mathrm{HF}}
    =
    \frac{1}{2}\sum_{pq} n_p n_q
    \left[
      (pp|qq) - (pq|qp)
    \right].
  \end{equation}
  \item Compare active-space exact diagonalization with the ADAPT-VQE result.
\end{itemize}

\section{Near-Term Software Tasks}

\begin{enumerate}[leftmargin=*]
  \item Keep H$_2$ as the regression test for integral construction.
  \item Add a TESTS case for isolated H$_2$CO.
  \item Add a TESTS case for H + H$_2$CO at one reactant-like geometry.
  \item Add a data object that stores $E_{\mathrm{nuc}}$, $h_{pq}$, $h_{pqrs}$, MO coefficients, occupations, and metadata.
  \item Add active-space slicing of $h_{pq}$ and $h_{pqrs}$.
  \item Add Jordan-Wigner mapping for the active-space molecular Hamiltonian.
  \item Validate against exact diagonalization before connecting to ADAPT-VQE.
\end{enumerate}

\section{First Publishable Arc}

The first paper-level story can be:
\begin{quote}
Adaptive quantum active-space simulations of hydrogen addition to formaldehyde, with gas-phase
benchmarks and first water-ice barrier shifts.
\end{quote}
The first results should include:
\begin{itemize}[leftmargin=*]
  \item gas-phase reaction profile for CH$_3$O and CH$_2$OH channels,
  \item active-space Hamiltonian construction and validation,
  \item exact active-space and ADAPT-VQE energies,
  \item first comparison between gas phase and one-water or small-cluster ice embedding,
  \item resource estimates as a function of active-space size.
\end{itemize}

\section{Immediate Next Step}

The next concrete system after H$_2$ is isolated formaldehyde, H$_2$CO. This should be used to test
the Psi4 integral pipeline for a real astrochemical molecule before adding the incoming H atom or ice.

\end{document}
